\scanpage{68}
\mspage{- 19 -}\hfill n\textsuperscript{o}\handwritten{357}

\uline{pour tout} $\handwritten{\lambda}\handwritten{\in}\handwritten{\mathscr{C}}'(M)$, $U(x)$ \uline{est scalairement essentiellement} $\handwritten{\lambda}$-\uline{intégrable au sens de la topologie de la convergence compacte dans} $\handwritten{\mathscr{L}}(E,E)$ ; \uline{on a} $U(\handwritten{\lambda})=\handwritten{\int} U d\handwritten{\lambda}\handwritten{\in}\handwritten{\mathscr{L}}(E,E)$, \uline{et}, \uline{quels que soient} $u\handwritten{\in} E$ \uline{et} $u'$ \uline{dans le dual} $E'$ \uline{de} $E$ :
\begingroup\makeatletter\tagsleft@true\makeatother
\[
U(\handwritten{\lambda})u=\handwritten{\int} U(x)u.d\handwritten{\lambda}(x)\quad,\quad
\handwritten{\langle} U(\handwritten{\lambda})u,u'\handwritten{\rangle}=\handwritten{\int}\handwritten{\langle} U(x)u,u'\handwritten{\rangle}\,d\handwritten{\lambda}(x)\quad;\tag*{(9)}
\]
\uline{de plus}, \uline{quels que soient} $\handwritten{\lambda}\handwritten{\in}\handwritten{\mathscr{C}}'(M)$, $\mu\handwritten{\in}\handwritten{\mathscr{C}}'(M)$, \uline{on a} :
\[
U(\handwritten{\lambda}\ast\mu)=U(\handwritten{\lambda})U(\mu)\ .\tag*{(10)}
\]
\endgroup
\uline{Si enfin l'ensemble des} $U(x)$, \uline{pour} $x\handwritten{\in} M$, \uline{est équicontinu}, \uline{les mêmes résultats subsistent lorsqu'on substitue} $\handwritten{\mathscr{M}}^{1}(M)$ \uline{à} $\handwritten{\mathscr{C}}'(M)$.

Le fait que $U(\handwritten{\lambda})$ est défini et est dans $\handwritten{\mathscr{L}}(E,E)$ est un cas particulier de la prop.16 du Chap.VI, \handwritten{\S}\,1,n\textsuperscript{o}7 ; la seconde des relations (9) résulte alors de la définition même de $\handwritten{\int} U d\handwritten{\lambda}$, puisque $\handwritten{\langle} Uu,u'\handwritten{\rangle}$ est une forme linéaire continue sur $\handwritten{\mathscr{L}}(E,E)$ (prop.11 de E.V.T., Chap.IV, \handwritten{\S}\,2, n\textsuperscript{o}9, joint au fait que la topologie de la convergence compacte est plus fine que celle de la convergence simple) [N.B. - Le rédacteur est tout fier d'avoir pu donner une raison aussi savante pour un fait aussi évident, mais n'a pas trouvé dans Bourbaki d'autre référence pour celui-ci]. La première suit de là en vertu des définitions. [N.B. - Les formules (9) auraient en réalité dû figurer au Chap.VI, \handwritten{\S}\,1, n\textsuperscript{o}7 !]. En substituant $\handwritten{\lambda}\ast\mu$ à $\handwritten{\lambda}$ dans la seconde formule (9), et appliquant la définition de $\handwritten{\lambda}\ast\mu$, on obtient :
\[
\handwritten{\langle} U(\handwritten{\lambda}\ast\mu)u,u'\handwritten{\rangle}
=\handwritten{\int}\handwritten{\langle} U(xy)u,u'\handwritten{\rangle}\,d\handwritten{\lambda}(x)d\mu(y)
=\handwritten{\int}\handwritten{\langle} U(y)u,{}^tU(x)u'\handwritten{\rangle}\,d\handwritten{\lambda}(x)d\mu(y)\ .
\]
Par hypothèse, $x\handwritten{\longrightarrow} U(x)u$ est continue, donc $\handwritten{\langle} U(xy)u,u'\handwritten{\rangle}$ est continue sur $M\times M$ ; comme l'intégrabilité essentielle d'une fonction continue entraîne son intégrabilité, on peut appliquer le théorème de Lebesgue-Fubini au deuxième membre, ou, ce qui revient au même au troisième membre de la relation ci-dessus ; cela donne, d'après la première relation (9) :
% archive sheet number: 67

